\documentclass[11pt]{article}
\usepackage[margin=1in]{geometry}
\usepackage{amsmath,amssymb}
\usepackage{booktabs}
\usepackage{hyperref}
\usepackage{siunitx}

\title{Replication Report:\\
Lagaris, Likas, Fotiadis (1998) --- Artificial Neural Networks for Solving ODEs and PDEs}
\author{Replication effort, PDE-100 set}
\date{2026}

\begin{document}
\maketitle

\begin{abstract}
We independently replicate the trial-solution artificial neural network (ANN)
method of Lagaris, Likas, and Fotiadis (\emph{IEEE TNN} 9(5):987--1000, 1998),
which is the direct precursor of modern physics-informed neural networks
(PINNs). Three of five worked examples are re-run: Problem 1 (first-order
ODE), Problem 3 (second-order Dirichlet BVP), and Problem 5 (2D Poisson on the
unit square). All three reproduce the paper's reported accuracy band to within
a factor of $\sim$2, and the paper's central interpolation-superiority claim
(C4) is confirmed at $68.6\times$ on an ODE analog. \textbf{Verdict:
REPLICATED.}
\end{abstract}

\section{Paper and scope}
Preprint: arXiv \texttt{physics/9705023 v1} (19 May 1997, 26 pp). We test three
of the five worked examples from Section 4 of the paper (P1, P3, P5) plus the
qualitative interpolation-superiority claim (C4) from the Section 4.2 / Table 1
discussion.

\section{Method (as implemented)}
The trial solution
\[
\Psi_t(x) = A(x) + F(x)\,N(x,p)
\]
combines a hand-coded function $A(x)$ that satisfies the initial/boundary
conditions with $F(x)$ that vanishes on the initial/boundary set, so that the
BCs are satisfied \emph{exactly} for any parameter vector $p$. The MLP
$N(x,p)$ has one hidden layer of 10 sigmoid units and a linear scalar output
(no output bias); weights are initialized $U(-1,1)$, torch seed $=0$. Training
minimizes the sum of squared residuals of the differential operator on the
paper's collocation grids using \texttt{torch.optim.LBFGS} with strong-Wolfe
line search (\texttt{max\_iter} $\in\{4000, 5000, 5000\}$, \texttt{history\_size}
$= 50$, tolerances \num{1e-14}). Language/packages: Python 3.12, PyTorch 2.2.2
(CPU, \texttt{float64}), NumPy 1.26.4. Deterministic under seed 0.

\subsection{Trial forms (verbatim from paper Sec. 4)}
\begin{align}
\text{P1:}\quad & \Psi_t(x) = 1 + x\,N(x,p) \\
\text{P3:}\quad & \Psi_t(x) = x\sin(1)e^{-1/5} + x(1-x)\,N(x,p) \\
\text{P5:}\quad & \Psi_t(x,y) = A(x,y) + x(1-x)y(1-y)\,N(x,y,p)
\end{align}
where $A(x,y)$ is the closed-form bilinear boundary lift of paper Eq.~18.

\section{Results}
\begin{table}[h]
\centering
\small
\begin{tabular}{lrrrrrl}
\toprule
Problem & Grid & Iters & Final loss & max err & RMSE & Paper \\
\midrule
P1 (1st-order ODE) & 10 tr / 200 te & 4760 & \num{4.10e-7} & \num{2.72e-5} & \num{1.51e-5} & Fig.~2: $\sim 10^{-6}$--$10^{-5}$ \\
P3 (2nd-order BVP) & 10 tr / 200 te &  594 & \num{3.17e-7} & \num{3.54e-6} & \num{1.97e-6} & Fig.~5: $\sim 10^{-6}$--$10^{-5}$ \\
P5 (2D Poisson)    & $10\times 10$ tr / $41\times 41$ te & 6250 & \num{2.71e-7} & \num{9.59e-7} & \num{3.47e-7} & Table~1: \num{5e-7} \\
\bottomrule
\end{tabular}
\caption{Per-problem replication results. All three within factor $\sim 2$ of paper.}
\end{table}

\subsection{C4 --- interpolation-superiority (Problem 1, same 10 training nodes)}
\begin{table}[h]
\centering
\begin{tabular}{lrr}
\toprule
Method & max err @ train nodes & max err @ 200-pt dense grid \\
\midrule
Trapezoid FD + natural cubic spline & \num{7.99e-4} & \num{1.86e-3} \\
ANN (this replication)              & \num{2.72e-5} & \num{2.72e-5} \\
\midrule
Ratio FD/ANN on dense grid          &               & $\mathbf{68.6\times}$ \\
\bottomrule
\end{tabular}
\caption{Off-grid generalization: ANN error uniform, FD error concentrated at nodes.}
\end{table}

The ANN's error is essentially uniform between training and dense-grid
evaluation, exactly as the paper claims for its PDE cases (Table~1: neural
\num{5e-7} at both training and $30\times 30$ test; FEM \num{2e-8} at training
but \num{1.5e-5} at test). Qualitative phenomenon reproduces.

\section{Verdict}
\textbf{REPLICATED.} LLM judge (Argo \texttt{gpt-5}) independently returned
\texttt{verdict = REPLICATED}, \texttt{coverage\_score = 10/10},
\texttt{agreement\_score = 9/10}.

\section{Genuine critique}
This section is a candid appraisal of what our replication does \emph{not}
settle and where the paper itself is thin or over-generalizes. It is
deliberately separated from the mechanical claims ledger.

\paragraph{1.\ Scope of replication is deliberately narrow.}
We reproduced 3 of the 5 worked examples (P1, P3, P5). Problem 2 (very similar
to P1), Problem 4 (coupled first-order ODE system), and Problems 6--8 (harder
2D PDEs with mixed BCs) were not attempted. All are structurally analogous to
what was tested and would almost certainly reproduce, but ``analogous'' is not
``verified.'' The strongest interpretation of our verdict is: \emph{the method
as described works on the three cases we tested at the reported accuracy}.

\paragraph{2.\ The C4 comparator is FD+spline, not FEM.}
The paper's PDE comparator in Table~1 is 2D FEM on $18\times 18$ elements. We
implemented an ODE analog (implicit trapezoid + natural cubic spline) because
we replicated the 1D case. The \emph{qualitative} phenomenon --- neural error
uniform, classical error concentrated at nodes and degrading off-grid ---
reproduces, and the ratio ($68.6\times$) is in the same order of magnitude as
the paper's FEM/ANN comparison. But we did not reproduce the specific FEM
numbers ($1.5\times 10^{-5}$ at test / $2\times 10^{-8}$ at training).

\paragraph{3.\ Optimizer is L-BFGS, not Merlin BFGS.}
Both are second-order quasi-Newton with line search. L-BFGS is the standard
modern substitute and it converges to the same accuracy regime. But the
paper's specific claim of ``1--5 iterations for linear problems'' refers to
Merlin's outer-iteration count in a particular way that is not directly
comparable to L-BFGS line-search-step counts (594--6250 for us). This is
\emph{not} a discrepancy; it is an apples-to-oranges accounting question that
neither we nor the paper resolves cleanly.

\paragraph{4.\ Single seed, no restart statistics.}
We report a single deterministic run at \texttt{torch.manual\_seed(0)}. P5's
max error ($9.59\times 10^{-7}$) is a factor of $\sim 2$ worse than the
paper's best figure ($5\times 10^{-7}$); a distribution over restarts would
likely straddle the paper's number, but we chose one honest run over cherry
picking. This is a strength for reproducibility and a weakness for
head-to-head comparison.

\paragraph{5.\ Method-level fragilities the paper does not stress.}
The trial-solution construction requires a hand-designed $A(x)$ that satisfies
the BCs and an $F(x)$ that vanishes on the boundary. In simple product
domains this is straightforward. In irregular domains, mixed BCs, or Neumann/
Robin BCs it is nontrivial and is the reason modern PINNs abandoned the
approach in favour of soft-loss boundary penalties. The paper does address
irregular domains in later sections (Problems 6--8, radial-basis boundary
lifts), but the difficulty grows quickly and is under-discussed relative to
the polish of the P1--P5 examples we replicated.

\paragraph{6.\ ``Interpolation superiority'' is a comparator-dependent claim.}
The paper's stronger claim --- ANNs interpolate better than FEM \emph{because}
$\Psi_t$ is a closed analytic form --- is only about the choice of ansatz, not
about any deep property of neural networks. A spectral method or an
appropriately-refined polynomial FEM (say, $p$-FEM) would also interpolate
smoothly off-grid. The paper's phrasing can be read as making a stronger
neural-specific claim than the experiments actually justify; the honest
statement is ``smooth global ansatz beats piecewise-local ansatz at
interpolation,'' which is a comparator artefact.

\paragraph{7.\ No stiffness, no nonlinearity, no high dimension.}
All five worked examples are smooth, low-dimensional, and (mostly) linear.
The paper does not test on stiff ODEs, strongly nonlinear PDEs, or dimensions
$>2$ --- exactly the regimes that motivated modern PINNs and where trial-
solution ANNs have not been shown to scale. Our replication inherits this
scope limitation.

\paragraph{8.\ Adaptive collocation is not addressed.}
All five worked examples use uniform collocation grids. Modern PINN practice
has moved toward residual-adaptive and importance-sampling collocation because
uniform grids waste effort in smooth regions. The paper does not discuss the
adaptivity question at all --- reasonable in 1998, but relevant to how one
should interpret its accuracy figures today.

\section{Honest gaps (from ledger)}
\begin{enumerate}
\item Only 3 of 5 worked examples reproduced.
\item C4 comparator is FD+spline, not FEM.
\item L-BFGS substituted for Merlin BFGS (standard substitution).
\item Single seed reported; P5 within $\sim 2\times$ of paper's best.
\item Text extraction used \texttt{pdftotext -layout} (OCR failed on the PDF);
      equations verified against arXiv.
\end{enumerate}

\section{Reproduction}
\begin{verbatim}
source work/venv/bin/activate
python3 work/lagaris_ann.py report/evidence   # ~30 s
python3 work/llm_judge.py  report/evidence    # ~30 s
\end{verbatim}
Deterministic under seed 0; L-BFGS-CPU double is deterministic.

\end{document}
